\section{极值分布与尾部风险}
\label{sec:05-Probability/06-Common-Distributions/08}

上一节盯着``平均而言会怎样''，以及这个问题的失效边界。但有一类决策从头到尾就不关心平均。

你被派去设计一段海堤。手上有过去 50 年每年最高水位的记录——3.2 米、2.8 米、4.1 米、3.5 米，等等。问题不是``平均每年水位多高''，你不需要挡住平均那一次。堤坝必须在使用期内挡住最高的那一次。如果设计水位定低了，淹一次就够你受的；定高了，每一厘米多花上千万。风险管理关心最大单日亏损；链条的承载能力取决于最弱一环。把对象从``平均''换成``最极端的一次''，极限理论几乎要重建一遍。

你在历史记录里翻到一个 4.7 米，那是 50 年的最大值。你觉得在堤坝上写``设计水位 4.7 米''够安全吗？直觉告诉你——五十年内的最大值不等于一百年内的最大值。那把 4.7 米乘以一个安全系数？系数的依据是什么？你发现你在用一个没经过思考的数字去保护一笔想清楚了的投资。

这时候你也许会想：能不能对极端值算一个置信区间？先把每年的最大值当成某种分布的样本，然后像中心极限定理对待样本均值那样，找一个极限分布来逼近最大值的分布。

好消息是，最大值也有类似中心极限定理的结构：存在普适的极限分布族。坏消息是，极限不再是 Gaussian——它归入另外三族。而上一节的尾指数，恰好决定你落入哪一族。

\subsection{块最大值本身没有极限分布}

给数据分块是一个自然的操作：年最大日降雨量，月最大单日亏损——每块取一个最大值。这块的大小有讲究。块要够大，大到里面的最大值摆脱了日常琐碎波动，只反映机制的尾部特征；但不能大到跨越不同机制，把来源不同的极端值搅进同一块。选年作为自然周期是一种方便，但不是数学要求。

设 \(X_1,\ldots,X_n\) 独立同分布，分布函数 \(F\)，记块最大值

\[
M_n = \max\{X_1,\ldots,X_n\}.
\]

它的分布函数有一个不用近似的精确表达式：

\[
P(M_n\le x) = P(X_1\le x,\ldots,X_n\le x) = F(x)^n.
\]

这个表达式干净，但直接使用它有问题。取一个具体数字感受一下：假设每年最高水位服从某个分布，上端无界。你盯着任一个固定的水位 \(x\)。因为 \(F(x) < 1\)，取 \(n\) 次方之后，\(F(x)^n \to 0\)——最大值一路向右漂移，离你选的 \(x\) 越来越远。如果你选的分布上端无界，最大值发散到无穷；如果上端有界，最大值退化到那个端点常数。无论哪种情况，直接看 \(M_n\) 都得不出有意义的极限分布。换句话说，块数和块内观测数一起增长时，最大值本身的分布不稳定。

你拿着 50 个年最大值去拟合一个分布，假如明年起你收集 100 个年最大值，这 100 个最大值的分布和那 50 个的分布不是同一族——它的位置往右跑了。一个随样本量漂移的分布没法拿来外推。你没法说``五十年数据的最大值分布的 99\% 分位数就是设计水位''，因为一百五十年数据的最大值分布压根不是同一个东西。

这个困境中心极限定理早就面对过。回想一下：\(\sum_i X_i\) 本身也不收敛，需要先中心化减去 \(n\mu\)，再除以 \(\sqrt{n}\)，才收敛到 Gaussian（见\hyperref[sec:05-Probability/08-Limit-Theorems/03]{``中心极限定理与 Gaussian 近似''}）。和量发散，但经过合适的平移和缩放后，有一个稳定的极限形状。对最大值，照搬这份智慧：不对 \(M_n\) 本身找极限，而是找规范化常数 \(a_n>0\) 和 \(b_n\)，问

\[
P\Bigl(\frac{M_n-b_n}{a_n}\le x\Bigr) = F(a_n x + b_n)^n
\]

能否收敛到一个非退化的分布 \(G\)。对象从和换成了最大值，平移和缩放的量也不同，但问题的结构一模一样：什么样的 \(G\) 可能出现？需要什么样的 \(F\) 才能让 \(G\) 存在？

\subsection{Fisher--Tippett--Gnedenko：只有三种归宿}

答案由 Fisher--Tippett--Gnedenko 定理给出（证明超出本单元范围）：若存在规范化常数使上式收敛到非退化的 \(G\)，则 \(G\) 必然属于广义极值分布（GEV）族：

\[
G(x) = \exp\Bigl\{-\bigl[1+\xi\bigl(\frac{x-\mu}{\sigma}\bigr)\bigr]^{-1/\xi}\Bigr\},
\]

定义域为 \(1+\xi(x-\mu)/\sigma>0\)。先别被这个公式吓到。它就是一个三参数分布族，参数各有直观含义：\(\mu\) 控制极值中心在哪（位置参数），\(\sigma>0\) 控制极值散布多宽（尺度参数），而第三个参数 \(\xi\)——形状参数——控制尾部到底有多厚，它把整个族劈成三个子类。\(\xi=0\) 的情况需要按 \(\xi\to0\) 取极限，得到 \(G(x)=\exp\{-e^{-(x-\mu)/\sigma}\}\)。

{\def\LTcaptype{none}
\begin{longtable}{@{}llp{8cm}@{}}
\toprule\noalign{}
极限族 & \(\xi\) & 母体尾部类型与代表 \\
\midrule\noalign{}
\endhead
\bottomrule\noalign{}
\endlastfoot
Gumbel & \(\xi=0\) & 轻尾、指数级衰减、上端无界：Gaussian、Exponential、Gamma、对数正态 \\
Fr\'echet & \(\xi>0\) & 幂律重尾，尾指数 \(\alpha=1/\xi\)：Pareto、Student \(t\)、Cauchy \\
Weibull & \(\xi<0\) & 有有限上端点：Uniform、Beta \\
\end{longtable}
}

三种归宿对应三种直觉画面，这不是数学分类，而是自然界的三种尾部地貌。

Gumbel 对应的母体尾巴按指数衰减。指数衰减很快——快到你取一百万个样本，最大值也只比典型值大几个标准差。画面是大自然用一堵指数墙压住了极端值：再大的 \(n\)，最大值也只能慢慢往上蹭。轻尾世界里，极端值被箍住了。

Fr\'echet 对应的母体尾巴按幂次衰减，速度慢得多。幂次尾上没有墙，只有越来越长的缓坡。最大值可以按 \(n^{1/\alpha}\) 的幂次冲出——如果 \(\alpha\) 很小（尾很重），几百个样本里的最大值就可以比典型值大几个数量级。取一百万个样本，那更是大到让你瞠目。

Weibull 对应母体有硬上限。考试满分 100 分，材料的理论强度有个物理上限——收集再多样本，最大值也只能逼近这道天花板，无法突破。

这里有一个容易被跳过的关键点：落入哪一族由母体分布的尾部性状决定，不是建模者可以自由选择的。你的数据来自一个幂律尾的过程，你偏要塞进 Gumbel——这跟在重尾数据上套 Gaussian 置信区间是同一种结构错误，只是换了一层包装。反过来，给有硬上限的数据套 Fr\'echet，你会外推出物理上不可能的超限值。把三族写进一个带连续形状参数 \(\xi\) 的 GEV 公式的实用价值正在于此：从数据估计出 \(\xi\) 的正负和大小，本身就是对母体尾部的一份诊断报告——它替数据告诉你，你面对的是墙、是坡，还是天花板。

\subsection{尾指数决定极限族，也决定最大值长多快}

上一节的尾指数在这里直接变现。若母体尾部 \(P(X>x)\sim x^{-\alpha}\)，规范化最大值收敛到 Fr\'echet 族，且 \(\xi=1/\alpha\)——\(\alpha\) 越小、尾越重，\(\xi\) 越大，极限分布自己的上尾也越厚。尾指数不只是描述母体的一个数字，它还锁定了极端值的极限行为。

最大值的典型量级同样由 \(\alpha\) 控制。拿 Pareto 母体，取规范化 \(a_n=x_m n^{1/\alpha}\)（无需平移，\(b_n=0\)）：

\[
P\Bigl(\frac{M_n}{x_m n^{1/\alpha}}\le x\Bigr)
= \Bigl(1-\frac{x^{-\alpha}}{n}\Bigr)^n
\to e^{-x^{-\alpha}},
\]

即最大值与 \(n^{1/\alpha}\) 同阶增长。注意这个增长速度直接由尾指数决定——尾指数减半，最大值的幂次翻倍。

作为对照，Gaussian 母体的 \(M_n\) 典型量级只有 \(\sqrt{2\log n}\)。这个 \(\sqrt{2\log n}\) 是怎么来的？思路是这样的：先确定位置参数 \(b_n\)，让每块平均恰好有一个观测超过它——即 \(nP(X>b_n)\approx 1\)。对 Gaussian 尾部用 Mill 近似 \(P(X>x)\approx e^{-x^2/2}/(x\sqrt{2\pi})\) 解这个方程，主导项就是 \(b_n\approx\sqrt{2\log n}\)。轻尾最大值的对数增长，根源正是尾部的指数级衰减——衰减太快了，取更多样本也只能挤出对数级别的额外极端值。

代入 \(n=10^6\) 让两种增长模式的差距落到数字上：

\begin{itemize}
\tightlist
\item Gaussian：最大值约 \(5\) 倍标准差。一百万观测的天文数字，里头最大的那个也只比均值多出 5 个标准差——连两位数都不到；
\item Pareto \(\alpha=2\)：最大值在 \(10^{6/2}=10^3\) 倍尺度参数的量级——三个数量级；
\item Pareto \(\alpha=1\)（Cauchy 尾部）：最大值在 \(10^6\) 倍尺度的量级——单个最大值和样本量 \(n\) 本身同阶。你再取一百万个样本，最大值的量级又涨一百万倍。
\end{itemize}

尺度参数和标准差不能直接换算出精确倍数，这个对比重要的不是具体数字，而是增长方式上的质变：轻尾世界里，最大值按对数缓慢爬行——你费劲收集更多数据，极端值只吝啬地涨一点点；重尾世界里，它按幂次冲出——每多收集一个数量级的样本，最大值就多蹦一个数量级。上一节说``平均无法稀释极端值''，这里看到的是同一事实的另一半——极端值主导的范围随 \(n\) 扩张的速度，决定了平均稀释策略何时彻底失败。如果你在 Pareto \(\alpha=1\) 的世界里靠增加样本量来压低最大值的相对影响，你会绝望地发现它跑得比你快。

\subsection{阈值法：不浪费每一条极端信息}

块最大值方法有一个挥霍信息的问题：每块只保留一个数，块里第二、第三大的值——同样携带尾部信息——被全部丢弃。按年分块，五十年观测只换来五十个最大值，拟合 GEV 的样本量少得可怜。五十个点去估计三个参数，还指望外推到百年一遇的水平——你手上那点信息本来就薄。

有没有办法不浪费那些次极端的观测？有。更经济的做法是阈值超限（POT，peaks over threshold）：固定一个高阈值 \(u\)，把所有超过 \(u\) 的观测都纳入分析，研究超出量 \(X-u\given X>u\) 的分布。你不用再把数据剁成年块然后每块只捡一个最大值——所有超过阈值的观测都是有效样本。五十年日水位数据里，超过某个高阈值的可能有两三百个观测，比五十个年最大值富裕得多。

Pickands--Balkema--de Haan 定理给出了与块最大值平行的结论：对足够高的 \(u\)，超出量近似服从广义 Pareto 分布（GPD）：

\[
P(X-u\le y\given X>u) \approx
1-\Bigl(1+\frac{\xi y}{\beta}\Bigr)^{-1/\xi},
\qquad y\ge 0,\; 1+\frac{\xi y}{\beta}>0,
\]

\(\beta>0\) 是取决于 \(u\) 的尺度参数，\(\xi=0\) 时按极限理解为 Exponential。形状参数 \(\xi\) 与 GEV 中的是同一个数字——母体尾指数再次登场。这不是巧合：块最大值和阈值法是同一极限规律的两种取样方式，它们共享同一个尾部形状参数。数据昂贵时，不浪费任何极端观测的阈值法通常更高效，代价是需要额外选择 \(u\)：太低则近似失效（普通波动混进来，GPD 不适用），太高则超阈值数据所剩无几（用来拟合的样本又变回个位数）。

选 \(u\) 有一个常用的诊断工具。当超出量确实服从 GPD 且 \(\xi<1\)，均值超出函数

\[
e(u) = \E[X-u\given X>u]
\]

是 \(u\) 的线性函数。操作上，把样本均值超出量对 \(u\) 作图——横轴是不断升高的阈值，纵轴是超过该阈值的观测的平均超出量——曲线开始呈现直线的起点就是候选阈值。在这个点之前，你把普通波动当极端处理，曲线乱晃；过了这个点，真正的尾部结构浮现，线性关系登场。这个诊断和上一节的 Hill 估计面临同样的困境：对选择敏感，收敛慢，需要大量极端数据——所有尾部诊断方法共享这条边界。如果在直线段开始的地方样本已经太少，你没法同时做到``阈值够高''和``样本量够拟合''。

\subsection{``百年一遇''到底在说什么}

极值分析的最终产出常常是重现水平。\(T\) 年重现水平 \(z_T\) 定义为每年有概率 \(1/T\) 被超越的量：

\[
P(M > z_T) = \frac{1}{T}.
\]

给定拟合好的 GEV，\(z_T\) 从 \(G(z_T)=1-1/T\) 解出。

``百年一遇洪水''的精确含义是``每年以 \(1\%\) 的概率被超越''。它不保证一百年内只发生一次。如果各年近似独立且概率稳定，一百年内至少遭遇一次的概率是

\[
1-\Bigl(1-\frac{1}{100}\Bigr)^{100} \approx 63\%,
\]

超过一半。把``百年一遇''理解为``一百年安全''，是把重现期和安全期混为一谈。一百年间至少来一次的概率是 63\%，不等于 100\%。更扎心的是，十年内至少遭遇一次``百年一遇''的概率也有 \(1-(0.99)^{10}\approx 9.6\%\)——不是零。

而且 \(z_T\) 是被估计出来的：\(T\) 越大，在数据范围之外外推越远，置信带越宽。报告里的``百年一遇''常常挂着数倍的误差范围——意味着真实的百年一遇水位可能在报告数字的一半到两倍之间。这不是估计方法不好，而是从有限数据外推到极端分位数时，不确定性爆炸式增长。

更深层的边界在于外推的假设。从五十年数据估计百年一遇的水平，已经踩到了数据范围之外——你在用五十年信息去猜一百年的事，靠的是极限定理的普适性。但极限定理成立的前提是平稳性：产生极端值的物理过程或制度过程，在观测期和外推期内保持同一个尾部。

气候变化正在移动洪水的尾部——过去五十年统计出的尾部和未来五十年的尾部可能已经不是同一个；金融制度和市场结构的变化会移动亏损的尾部——2008 年前的尾部参数和 2008 年后的尾部参数活在两个世界里。历史尾部一旦失真，再精细的 GEV 拟合也只是精确地描述了一个不再存在的世界。这不是数学的失败——数学已经给了你最佳的平稳外推——而是外推这件事本身的边界：你只能假设明天跟昨天长得像，但这个假设不归数学管。

此外，块最大值理论默认观测近似独立。极端事件若成簇到来——暴风雨连续数日、市场恐慌持续数周——独立性假设失效，需要引入 extremal index 等修正，这超出了本单元范围。簇状极端事件的年最大值不会独立，因为同一年内的多次极端值可能来自同一场持续事件。

极值理论把``最坏会怎样''从修辞变成了可计算的分布问题。它的价值不在于给出一个精确数字——你知道那个数字后面挂着多宽的置信区间——而在于把尾部风险放进一个可以讨论和比较的框架：你可以说``按目前的尾部估计，百年一遇的水位约在 a 米到 b 米之间''，而不是只能说``大概挺大的''。它的边界同样清晰：所有定量结论都系在尾部平稳和足够的极端数据上。外推得越远，缆绳越细。在数据范围的两倍、三倍之外——比如从五十年数据外推到五百年一遇——你已经在缆绳将断未断的区域行走，那时更需要的是判断``这个外推还合理吗''，而不是算那个数字的小数点后几位。
